% Generated draft: requires mathematical and source review.
\documentclass[11pt]{article}
\usepackage[margin=25mm]{geometry}
\usepackage{fontspec}
\setmainfont{Latin Modern Roman}
\usepackage{amsmath,amssymb,mathtools,unicode-math}
\setmathfont{Latin Modern Math}
\usepackage{enumitem,needspace,xurl,hyperref}
\setlength{\parindent}{0pt}
\setlength{\parskip}{5pt}
\setlength{\emergencystretch}{3em}
\newcommand{\N}{\mathbb{N}}
\newcommand{\Z}{\mathbb{Z}}
\newcommand{\Q}{\mathbb{Q}}
\newcommand{\R}{\mathbb{R}}
\newcommand{\C}{\mathbb{C}}
\newcommand{\F}{\mathbb{F}}
\newcommand{\A}{\mathbb{A}}
\newcommand{\PP}{\mathbb{P}}
\newcommand{\E}{\mathbb{E}}
\newcommand{\eps}{\varepsilon}
\newcommand{\dd}{\,\mathrm{d}}
\begin{document}
\section*{UnsolvedMath ID: 30002086}
\textbf{Source problem code:} catalogue entry
\subsection{Definitions and mathematical statement}
A gradient shrinking Ricci soliton is a triple \((M,g,f)\), where \((M,g)\) is a Riemannian manifold and \(f\in C^\infty(M)\) satisfies
\[
\operatorname{Ric}_g+\nabla_g^2 f=\lambda g
\qquad\text{for some constant }\lambda>0.
\]
Here \(\operatorname{Ric}_g\) is the Ricci tensor, \(\nabla_g^2f\) is the Hessian, \(\nabla_g f\) is the gradient, and \(d_g\) denotes Riemannian distance. Keeping \(\lambda>0\) unspecified avoids imposing a normalization absent from the supplied record.

For every complete gradient shrinking Ricci soliton \((M,g,f)\) and every chosen basepoint \(p\in M\) satisfying
\[
f(p)=\min_{y\in M}f(y),
\]
is there a constant \(C=C(M,g,f,p)>0\) such that
\[
\forall x\in M,\qquad
 d_g(x,p)\geq C
 \quad\Longrightarrow\quad
 |\nabla_g f|_g(x)\geq C^{-1}?
\]
The original wording asks for a compact subset \(K\subseteq M\) and a constant \(c>0\) such that \(|\nabla_g f|_g(x)\geq c\) for every \(x\in M\setminus K\). The displayed distance formulation is the supplied cleaned statement. No noncompactness assumption is supplied; on compact \(M\), the distance condition can have an empty range of applicability.

\emph{Unresolved dimensional scope.} The supplied statements do not restrict \(\dim M\), whereas the background and literature assessment describe a condition for four-dimensional shrinkers. Thus the intended quantification may be restricted to \(\dim M=4\), or may range over arbitrary dimensions. The supplied evidence does not settle this distinction; the displayed question records the unrestricted wording without claiming that this broader scope was intended in the cited report.
\subsection{Short English statement}
Must the gradient of the potential of a complete shrinking Ricci soliton stay uniformly away from zero sufficiently far from a point where the potential attains its minimum? The bound and distance threshold may depend on the soliton and basepoint. The supplied context emphasizes dimension four, but the stated question leaves the dimension unspecified.
\subsection{Sources}
\begin{itemize}
\item Supplied local UnsolvedMath registry record: permanent ID \(30002086\), title \emph{Gradient Lower Bounds for Shrinking Ricci Solitons}, source problem code \texttt{catalogue entry}, in the Oberwolfach Reports Open Problems collection. No individual proposer is supplied; the proposed year is recorded as 2012.
\item Supplied report citation: \emph{Geometrie (2013)}, Oberwolfach Reports, DOI \texttt{10.4171/owr/2012/27}, pp.\ 1639--1686. Supplied source URL and literature-source link labelled ``original Oberwolfach report'': \url{https://doi.org/10.4171/owr/2012/27}. The title year is retained as supplied, alongside the record's 2012 proposed year and DOI identifier.
\item Supplied literature source: Haslhofer, \emph{Singularities in 4d Ricci flow}, \url{https://www.math.utoronto.ca/roberth/papers/Haslhofer_owr_singularities.pdf}. This attribution identifies the supplied reference, not an explicitly identified proposer of the problem.
\item Literature assessment checked at \texttt{2026-08-21}: the supplied curation labels the problem open and reports that no general proof of the technical condition for four-dimensional Ricci shrinkers was found. This is supplied dated curation, not an independently verified literature assessment. The background explicitly calls for independent verification of time-sensitive status and cited claims.
\item The registry labels the statement \texttt{corrected\_verified} and supplies the verification claim that the official report displays this condition and says that proving it for every shrinker remains open. These are supplied registry assertions; neither the cited report nor those assertions was independently checked for this draft.
\end{itemize}
\end{document}
